Un objecte es mou en línia recta i la posició, en metres, ve donada per
\[
s(t)=t^2+2t ,
\]
on $t$ és el temps en segons.

\begin{apartats}

\begin{tria}{tvm-tasca}
\itemtria{original}{0,75}{1,25}
Calcula la velocitat mitjana entre $t=1$ i $t=4$, i entre $t=1$ i $t=2$.

\begin{solucio}
$s(1)=3$, $s(2)=8$ i $s(4)=24$.\\
Entre $t=1$ i $t=4$: $\dfrac{24-3}{3}=7$ m/s. \quad Entre $t=1$ i $t=2$:
$\dfrac{8-3}{1}=5$ m/s.\\
La velocitat mitjana és la taxa de variació mitjana de la posició.
\end{solucio}

\itemtria{aproximacio-numerica}{0,75}{1,25}
Calcula la velocitat mitjana entre $t=1$ i $t=1{,}1$, entre $t=1$ i $t=1{,}01$, i entre $t=1$ i
$t=1{,}001$. Cap a quin valor sembla que s'acosten? Què representa aquest valor?

\begin{solucio}
$v_m=\dfrac{s(1+\Delta t)-s(1)}{\Delta t}$, amb $s(1)=3$:
\[
\frac{s(1{,}1)-3}{0{,}1}=\frac{3{,}41-3}{0{,}1}=4{,}1,\qquad
\frac{s(1{,}01)-3}{0{,}01}=4{,}01,\qquad
\frac{s(1{,}001)-3}{0{,}001}=4{,}001 .
\]
S'acosten a $4$ m/s: és la \textbf{velocitat instantània} en $t=1$, el límit de les velocitats
mitjanes quan l'interval es fa cada vegada més petit.
\end{solucio}
\end{tria}

\apartat[1,25]{1}
Calcula la velocitat instantània en $t=1$ aplicant la definició de derivada.

\begin{solucio}
\[
s'(1)=\lim_{h\to0}\frac{s(1+h)-s(1)}{h}
=\lim_{h\to0}\frac{(1+h)^2+2(1+h)-3}{h}
=\lim_{h\to0}\frac{h^2+4h}{h}=\lim_{h\to0}(h+4)=4 \ \text{m/s}.
\]
\end{solucio}

\begin{nomesllarg}
\apartat{0,75}
Calcula la velocitat instantània en $t=4$ i compara-la amb la velocitat mitjana entre $t=1$ i
$t=4$.

\begin{solucio}
\[
s'(4)=\lim_{h\to0}\frac{(4+h)^2+2(4+h)-24}{h}=\lim_{h\to0}\frac{h^2+10h}{h}=10 \ \text{m/s}.
\]
La velocitat mitjana entre $t=1$ i $t=4$ era de $7$ m/s, just entre els $4$ m/s de l'inici i
els $10$ m/s del final: l'objecte va accelerant.
\end{solucio}
\end{nomesllarg}

\end{apartats}
